\cotefile{156-9}
\scanpage{26}
\mspage{22}
% archive sheet number: 25 (upper left)

est \unsure{indispensable}, si on veut pouvoir \unsure{composer|comparer} tous les \uline{\unsure{intervalles}} d'une figure complète $\unsure{\Phi_T|\varphi_T}$ $(T\in\operatorname{Prop}(J))$ telles que
\[
\{J_{<a},\ J_{>a}\}\qquad\text{si }a\in J
\]
ou
\[
J_{<a},\ J_{a,b},\ J_{>b}\qquad\text{si }a<b,\ a,b\in J.
\]

\deleted{Dans} Le cas où $I=\{\omega_0,\omega_1\}$ correspond au cas où $J=\varnothing$, qu'il ne faut nullement exclure ! On trouve alors $\unsure{M|\mathcal M}$ réduit au seul élément $[J]$. On trouve la maquette réduite à une seule \unsure{str.t|Ar.t|Art}, mais visualisée comme figurant un \uline{intervalle ouvert indivisible} : l'intervalle ouvert indivisible est représenté par l'ens ordonné vide !

Si on veut une \unsure{formulation générale} commune, il faut dire qu'on se donne
\[
(I,\Omega_0;\Omega_1)
\]
avec
\[
I,\ \text{ens ordonné},\qquad\Omega_0,\Omega_1\subset I
\]
et les conditions
\margin{donc $\operatorname{Card}\Omega_0,\operatorname{Card}\Omega_1,\leq1$ et il y a quatre cas à distinguer}
\[
\left\{
\begin{aligned}
\text{a)}\quad&\deleted{\unsure{x}}\ x\in\Omega_0\deleted{\unsure{x}}\Longrightarrow x\text{ est un plus petit él. de }I\\
\text{a')}\quad&y\in\Omega_1\Longrightarrow y\ \text{\rule{2.7cm}{0.35pt}\ grand \rule{2.5cm}{0.35pt}}\\
\text{b)}\quad&x\in\Omega_0,\ y\in\Omega_1\Longrightarrow x\neq y.
\end{aligned}
\right.
\]
